\chapter{Literature Review}
\label{ch:literature-review}

\section{Introduction}
\label{sec:lr-intro}

This chapter reviews the published literature relevant to the four pillars of NirmalNode: (i) the scale and health impact of industrial and ambient air pollution, with particular attention to Bangladesh; (ii) Internet of Things (IoT)-based air quality monitoring systems; (iii) artificial intelligence (AI)/machine learning (ML)-based air pollution prediction; and (iv) air purification technologies, specifically HEPA filtration, activated-carbon adsorption, and cyclonic pre-separation. The chapter also reviews the emerging concepts of Green IoT, Green AI, and incremental/online (Adaptive) machine learning, which together form the methodological backbone of the proposed system. Each section closes with a short synthesis, and the chapter concludes with a consolidated comparison table and an explicit statement of the research gap that NirmalNode is designed to address.

\section{Industrial and Ambient Air Pollution: Global and Bangladesh Context}
\label{sec:lr-pollution-context}

Air pollution research has, over the last decade, converged on the finding that fine particulate matter (PM2.5) exposure carries health risks even at concentrations once considered safe. The World Health Organization's 2021 guideline revision reduced the recommended annual PM2.5 limit from 10 to 5 $\mu$g/m$^3$, and subsequent modelling work has shown that in many regions the \emph{non-anthropogenic} (natural) background component of PM2.5 already accounts for a substantial share of this reduced allowance, meaning that anthropogenic sources -- including industrial emissions -- must be reduced even further to remain within the new guideline \citep{pai2022who}. A parallel analysis focused on small and medium-sized cities concludes that the tightened WHO guideline is likely to be extremely difficult to meet for such cities without dedicated local-scale intervention, since they typically lack both the monitoring infrastructure and the regulatory capacity available to large metropolitan centres \citep{polezer2023new}.

For Bangladesh specifically, a national evidence review identifies vehicular emissions, brick kilns, and unplanned industrial growth as the principal drivers of the country's air pollution burden, and documents its substantial economic and public-health costs \citep{khatun2023evidence}. City-level and district-level studies corroborate this at finer resolution. Long-term reconstruction of PM2.5 exposure across Bangladesh using satellite and ground-based data identifies specific, persistent pollution hotspots rather than a spatially uniform exposure pattern, and links these hotspots to elevated long-term health risk \citep{ali2025longterm}. Within Dhaka, the Tejgaon industrial area shows PM2.5 and PM10 concentrations, and derived health-risk indices, that vary sharply with proximity to specific industrial sources \citep{basak2025spatial}, while satellite-based (Sentinel-5 TROPOMI) trend analysis of Chattogram -- the country's second-largest industrial hub -- shows that this city remains comparatively understudied relative to Dhaka despite carrying a substantial pollutant load \citep{huraira2024spatial}. Landfill-adjacent industrial zones in Dhaka have similarly been shown to emit and disperse gaseous pollutants (CH$_4$, CO$_2$, and other non-methane organic compounds) whose atmospheric impact had not previously been quantified \citep{aurnab2026monitoring}, and remote-sensing analysis of atmospheric CO$_2$ anomalies over Bangladesh links elevated concentrations to anthropogenic activity concentrated in and around industrial and urban centres \citep{ghosh2026spatiotemporal}. Industrial effluent near Dhaka has also been shown to elevate trace-element concentrations in the surrounding environment to levels of ecological and health concern \citep{lisa2025investigative}, and general population studies in pollution-prone urban areas of Bangladesh report measurable self-reported health deterioration associated with perceived air pollution exposure \citep{rahman2025perceived}.

\textbf{Synthesis.} This body of work establishes two facts that directly motivate the present thesis: (1) Bangladesh's air pollution problem is dominated by identifiable, spatially concentrated industrial and urban hotspots rather than being spatially uniform, and (2) existing evidence on these hotspots is drawn almost entirely from passive environmental monitoring (satellite remote sensing, periodic ground sampling) rather than from continuous, real-time, hotspot-embedded sensing capable of driving an automated response.

\section{Occupational Health Impacts at Industrial Hotspots}
\label{sec:lr-occupational-health}

A separate strand of literature documents the direct occupational health consequences of the pollutant exposures targeted by NirmalNode. Workers in fiber-cement roofing manufacturing exposed to elevated PM2.5 show measurable impairment of lung function \citep{nasri2023pm25}. Construction-site workers exposed to particulate matter show a range of associated adverse health effects \citep{sekhavati2023particulate}. Reviews of PM2.5 exposure across multiple industry types report associations with mental health conditions, including elevated risk of depression, anxiety, and, in some studies, more severe psychiatric outcomes, in addition to the well-documented respiratory and cardiovascular effects \citep{alhadhrami2024pm25mental}. These findings collectively indicate that the health burden of industrial air pollution is not limited to respiratory disease but extends to broader physical and psychological wellbeing, reinforcing the urgency of protecting workers at the specific points where their exposure is highest -- precisely the hotspots NirmalNode targets.

\section{IoT-Based Air Quality Monitoring Systems}
\label{sec:lr-iot-monitoring}

A substantial body of engineering literature demonstrates that low-cost microcontrollers and commodity sensors can be assembled into functioning air-quality monitoring networks. \citet{alam2025iot} describe an IoT-based real-time environmental monitoring system aimed specifically at developing-area deployment, emphasizing low hardware cost and remote accessibility. \citet{zhang2020realtime} combine mobile and fixed IoT sensing nodes to achieve localized air quality monitoring and prediction, demonstrating that a hybrid fixed/mobile sensing network can improve spatial resolution over fixed-only networks. \citet{kataria2022aiiot} propose a hybrid AI-and-IoT model for air quality prediction in a smart-city setting with network assistance, integrating communication-network parameters into the prediction pipeline. \citet{mani2021aipowered} present an AI-powered IoT system for real-time air pollution monitoring and forecasting, following a broadly similar sense-transmit-predict pipeline. At the systems level, \citet{ramadhani2025iot} demonstrate a low-cost-sensor IoT implementation for air quality monitoring, and \citet{pavan2025realtime} describe a comparable low-cost sensor-based IoT system with buzzer and SMS-based alerting and ThingSpeak-based remote visualization. More broadly, a review of low-cost sensor networks for air quality management in developing countries argues that such networks are not merely a cheaper substitute for reference-grade monitoring but a practically necessary tool given the financial and infrastructural constraints of these settings \citep{shabbir2025paradigm}. Complementary work applying AI to indoor air quality dynamics in developing nations similarly finds that field-deployed, low-cost sensing captures pollution patterns -- driven by cooking methods, ventilation, and living density -- that are highly context-specific and not well captured by generic models trained elsewhere \citep{karmakar2024exploring}. A review chapter on AI-based air pollution monitoring in developing countries reaches a parallel conclusion at policy scale, identifying data integration, infrastructure constraints, and energy-sector limitations as recurring barriers to AI deployment in this context \citep{noman2025aiair}. At a broader policy and enforcement scale, satellite analytics combined with AI and IoT sensing has also been proposed as a tool for regulatory enforcement against industrial air pollution in the South Asian context \citep{ahmed2025smart}.

\textbf{Synthesis.} These systems collectively establish IoT-based monitoring, and even AI-assisted prediction, as mature and well-validated concepts. However, in every reviewed case, the system's role ends at reporting or forecasting a pollutant level -- none of the reviewed IoT/AI monitoring systems couple their output to an automated, physical, localized mitigation action such as a purifier.

\section{Air Purification Technologies}
\label{sec:lr-purification}

\subsection{HEPA Filtration}
High-Efficiency Particulate Air (HEPA) filtration is the most extensively validated technology for removing fine particulate matter from air streams, with a defining performance standard of removing at least 99.97\% of particles of 0.3 $\mu$m diameter. Field evaluations confirm this effectiveness in real-world conditions: portable HEPA air cleaners have been shown to substantially reduce particulate concentrations in homes located close to major highways \citep{cox2018effectiveness}, and a controlled real-world evaluation of portable air cleaners confirms meaningful reductions in home particulate matter concentrations under everyday operating conditions \citep{lu2023realworld}. HEPA-equipped purification systems installed in intercity buses have been shown to reduce both particulate matter and airborne bacterial concentration \citep{lee2022effect}, and dedicated HEPA air cleaners have been shown to measurably improve indoor PM2.5 concentrations in enclosed residential and clinical settings \citep{chen2022efficacy}. Even under extreme pollution loading -- such as during wildland-urban-interface wildfire smoke intrusion -- portable HEPA filtration has been shown to meaningfully reduce indoor PM2.5 exposure \citep{chen2025fine}. A comparative assessment of HEPA-based air purifiers against ionization-based and other purification technologies concludes that mechanical HEPA filtration is more reliably effective, particularly for larger particle-size fractions, than ion-generator-based approaches \citep{dubey2021assessing}. A broader review of indoor air purifier technologies for microbial contamination reduction similarly finds HEPA-based mechanical filtration to be among the most consistently effective approaches evaluated \citep{szczotko2022review}, and dedicated microbiological-efficiency studies of HEPA filtration in medical-facility contexts report similarly favourable performance for airborne biological contaminants \citep{laine2025microbiological}.

\subsection{Activated Carbon Filtration}
While HEPA filters are effective against particulate matter, they are not designed to remove gaseous pollutants such as VOCs and odour-causing compounds. Activated carbon filtration, which relies on the adsorption of gas molecules onto the extensive internal pore surface area of activated carbon, is the standard complementary technology for this purpose and is a key stage in the purification train proposed for NirmalNode.

\subsection{Cyclone Separation}
Cyclonic (centrifugal) separation is a long-established pre-filtration technique that removes coarse particulate matter from an air stream by inducing tangential, swirling flow, throwing heavier particles outward against the separator wall under centrifugal force before the air stream reaches finer filtration stages. Placing a cyclone separator ahead of a HEPA filter, as proposed for NirmalNode, extends the operational lifetime of the (more expensive) HEPA stage by removing the coarse particulate load before it reaches the filter media.

\subsection{Adoption Barriers}
Despite the demonstrated effectiveness of these technologies, adoption in Bangladesh specifically remains very low: a multi-phase field study of household air purifier adoption in Dhaka finds that fewer than 1\% of middle-class households that can afford a purifier actually own one, driven substantially by misbeliefs about relative indoor/outdoor air quality and about purifier effectiveness, resulting in a willingness to pay far below the retail cost of a unit \citep{chowdhury2025misbeliefs}.

\textbf{Synthesis.} The purification literature confirms that HEPA and activated-carbon filtration, aided by cyclonic pre-separation, are well-validated technologies for both particulate and gaseous pollutant removal. However, every reviewed purification study evaluates these technologies as \emph{room-scale or vehicle-scale} appliances operating reactively on current air conditions -- none evaluate a compact, worker-proximate, hotspot-targeted unit operating on a predictive (1--2 hour look-ahead) activation schedule, which is the specific configuration proposed in this thesis.

\section{Green IoT and Green AI}
\label{sec:lr-green}

Green IoT refers to the design of IoT hardware and communication systems that minimize energy consumption, material use, and computational overhead while maintaining functional performance -- for example, through low-power microcontrollers, duty-cycled sensing, and efficient wireless protocols. Green AI extends the same philosophy to the machine learning layer: rather than defaulting to large, computationally expensive deep-learning models, Green AI favours lightweight models whose accuracy is sufficient for the task at hand while substantially reducing training time, inference latency, memory footprint, and energy draw. Within the reviewed air-quality literature, this philosophy is implicit rather than explicit: systems such as those of \citet{alam2025iot} and \citet{shabbir2025paradigm} are motivated primarily by cost constraints in developing-country deployment, which aligns closely with, but is not formally framed as, the Green IoT/Green AI paradigm that this thesis adopts explicitly as a design principle.

\section{Adaptive AI and Incremental (Online) Learning}
\label{sec:lr-adaptive-ai}

None of the AI-based air-quality prediction systems reviewed in Section~\ref{sec:lr-iot-monitoring} report the use of incremental or online learning; each trains a model once, on a fixed historical dataset, and deploys it statically. This is a significant limitation for any system intended for mass deployment across structurally different sites, since a model trained on one location's pollution signature (traffic patterns, industrial mix, meteorology) cannot be expected to generalize perfectly to a different site without some form of retraining or adaptation. Incremental (online) learning algorithms -- such as SGD Regressor, Passive-Aggressive Regressor, Hoeffding Trees, and Adaptive Random Forests, several of which are implemented in the open-source River library -- update model parameters continuously as new labelled data arrives, without requiring the full dataset to be reprocessed. This class of algorithm is well established in the broader online-learning literature but, to the best of the author's review of the supplied reference set, has not previously been applied to industrial, hotspot-level air pollution prediction coupled to an automated purification response. This absence constitutes the central methodological gap addressed by NirmalNode's Adaptive AI design, detailed in Chapter 3.

\section{Comparative Analysis of Reviewed Systems}
\label{sec:lr-comparison}

Table~\ref{tab:lr-comparison} summarizes the reviewed IoT/AI-based air-quality systems along the dimensions most relevant to NirmalNode: sensing scope, prediction capability, purification capability, and adaptivity of the underlying model.

\begin{table}[htbp]
\centering
\caption{Comparative summary of reviewed IoT/AI-based air quality systems}
\label{tab:lr-comparison}
\small
\begin{tabular}{@{}p{2.6cm}p{2.6cm}p{2.1cm}p{2.6cm}p{2.6cm}@{}}
\toprule
\textbf{System / Study} & \textbf{Sensing Scope} & \textbf{Predictive?} & \textbf{Purification Coupled?} & \textbf{Model Adaptivity} \\
\midrule
\citet{alam2025iot} & Fixed IoT node, developing-area & No & No & Static \\
\citet{zhang2020realtime} & Mobile + fixed IoT network & Yes (short-term) & No & Static \\
\citet{kataria2022aiiot} & City-scale, network-assisted & Yes & No & Static \\
\citet{mani2021aipowered} & Fixed IoT node & Yes & No & Static \\
\citet{ramadhani2025iot} & Fixed IoT node, low-cost & No & No & Static \\
\citet{pavan2025realtime} & Fixed IoT node, alert-based & No & No & Static \\
\citet{shabbir2025paradigm} (review) & Low-cost sensor networks & Mixed & No & Static \\
\citet{dubey2021assessing} & Room-scale purifier study & No & Yes (reactive) & N/A \\
\citet{lu2023realworld} & Home-scale purifier study & No & Yes (reactive) & N/A \\
\citet{lee2022effect} & Vehicle-scale (bus) purifier & No & Yes (reactive) & N/A \\
\citet{chen2025fine} & Home-scale, wildfire event & No & Yes (reactive) & N/A \\
\textbf{NirmalNode (proposed)} & \textbf{Localized hotspot node} & \textbf{Yes (1--2 hr ahead)} & \textbf{Yes (predictive)} & \textbf{Incremental / Adaptive} \\
\bottomrule
\end{tabular}
\end{table}

\section{Research Gap}
\label{sec:lr-research-gap}

Table~\ref{tab:lr-gap} consolidates the specific gaps identified across the preceding sections into the concrete design requirements that motivate NirmalNode.

\begin{table}[htbp]
\centering
\caption{Identified research gaps and corresponding NirmalNode design response}
\label{tab:lr-gap}
\small
\begin{tabular}{@{}p{4.3cm}p{4.0cm}p{4.6cm}@{}}
\toprule
\textbf{Gap Identified in Literature} & \textbf{Representative Sources} & \textbf{NirmalNode Design Response} \\
\midrule
Monitoring/prediction not coupled to an automated physical mitigation action & \citep{alam2025iot,zhang2020realtime,kataria2022aiiot,mani2021aipowered} & Prediction output directly drives purifier activation (Decision Engine, Ch. 3) \\
Purification evaluated only at room/vehicle scale, not worker-proximate hotspot scale & \citep{dubey2021assessing,lu2023realworld,lee2022effect,chen2022efficacy,chen2025fine} & Compact cyclone--HEPA--activated-carbon unit sized for a single hotspot / worker zone \\
Purifier activation is reactive (responds to current conditions only) & \citep{dubey2021assessing,lu2023realworld,lee2022effect} & 1--2 hour ahead predictive activation before hazard threshold is crossed \\
AI models are trained once and deployed statically; no post-deployment adaptation & \citep{alam2025iot,zhang2020realtime,kataria2022aiiot,mani2021aipowered} & Incremental/online learning (SGD, PA, Hoeffding Tree, ARF) for continuous, area-specific model update \\
Health/exposure evidence establishes hotspot severity but proposes no engineering countermeasure & \citep{nasri2023pm25,sekhavati2023particulate,alhadhrami2024pm25mental,rahman2025perceived} & Hotspot-targeted purification directly engineered around documented worker exposure points \\
Low-cost sensor deployment is necessary but insufficient without an integrated mitigation layer & \citep{shabbir2025paradigm,karmakar2024exploring,noman2025aiair} & Full sense--predict--purify pipeline built entirely on low-cost, locally available components \\
Very low real-world adoption of purification technology in Bangladesh due to cost and misbeliefs & \citep{chowdhury2025misbeliefs} & Low unit cost, localized (not whole-room) purification volume, and demonstrable predictive operation intended to reduce both cost and trust barriers \\
\bottomrule
\end{tabular}
\end{table}

\section{Chapter Summary}
\label{sec:lr-summary}

This chapter has reviewed four intersecting bodies of literature -- ambient and industrial air pollution in Bangladesh, IoT-based air quality monitoring, AI-based pollution prediction, and air purification technology -- and has shown that while each individual component of NirmalNode (low-cost IoT sensing, AI-based prediction, HEPA/activated-carbon/cyclone purification) is independently well validated in prior work, no reviewed system integrates all of them into a single hotspot-targeted, predictively-activated, incrementally-learning purification unit. This integration, and specifically the incorporation of incremental/online learning for area-specific adaptation, constitutes the central novelty of the system presented in the remainder of this thesis. Chapter 3 now presents the complete methodology by which NirmalNode realizes this integration.
